\section{Conclusion}
\label{sec:conclusion}
We asked whether a vessel forecast helps a nodule candidate scorer because it is placed correctly or merely because it adds information. A paired design in which control and lesion crops share their exterior, together with a value-preserving rotation sham, answers this question directly: on held-out data, alignment changes the location the scorer returns far more than the decision it makes, and a plain fixed subtraction is enough to realize this benefit. The same localization gain appears on scans with expert-annotated native vessels, where real anatomy rather than simulated tubes supplies the competing response.

A vessel forecast can appear useful for several reasons, since its shape and amplitude correlate with the type of background in a crop; combining a value-preserving rotation sham, an exterior-summary control and shared readouts is what isolates the contribution of its \emph{alignment}. Two practical implications follow. First, fixed subtraction combined with a compact-template search exploits the alignment as effectively as joint amplitude estimation, because a gain fitted inside the search region absorbs part of the lesion (Proposition~\ref{prop:absorb}), so it is a competitive baseline for more flexible nuisance models. Second, for pipelines that present a returned location to a reader, moving that location onto the lesion is valuable even when the number of accepted crops barely changes, whereas a gain in crop sensitivity can overstate the benefit when the returned location remains on the vessel.

The experiments also separate ranking from calibration. A matched-FPR effect establishes ranking separation at a given budget, whereas a shift in the negative score distribution changes the realized budget at an imported threshold. In our experiments, ranking gains from alignment transferred across acquisitions more reliably than calibrated cutoffs, so a deployment-oriented claim would require negatives from the target acquisition. A reliability-based rule for using the forecast selectively did not outperform always using it (Sec.~S6).

The study also has clear limits. All positive targets are inserted signals, so real nodules, scan-level proposal pipelines, partial-volume effects and per-scan false-positive rates remain untested. The prespecified crop-level criterion was not met, the localization robustness checks were conducted post hoc, and absolute localized sensitivity remains modest. Calibrated cutoffs, moreover, did not transfer across acquisitions even where ranking did.

These findings point to several directions. The most direct is a prespecified evaluation on adjudicated real nodules at the scan level with localized sensitivity as the primary endpoint. Because the learned CNN in our comparison used the forecast without exploiting its placement, a second direction is to build learned scorers that subtract or condition on an aligned forecast explicitly, which our results suggest is where the benefit originates. Replacing the straight-chord forecaster with one that follows curved and branching vessels should strengthen the effect on native anatomy, and calibrating thresholds with negatives from the target acquisition would turn the ranking gain into a transferable operating point. More broadly, a value-preserving placement sham is a simple and general control for any method that claims to benefit from spatial context.

